% F1 · La escala de identificabilidad (D6·02)
\documentclass[border=10pt]{standalone}
% ------------------------------------------------------------------
% Preámbulo común de las figuras del curso C1
% Paleta y tipografía alineadas con ../../estilo.tex
% ------------------------------------------------------------------
\usepackage[utf8]{inputenc}
\usepackage[T1]{fontenc}
\usepackage[spanish,es-noquoting]{babel}
\usepackage{helvet}
\renewcommand{\familydefault}{\sfdefault}
\usepackage{xcolor}
\usepackage{tikz}
\usetikzlibrary{arrows.meta,positioning,decorations.pathreplacing,calc,shapes.geometric,fit,backgrounds}

\definecolor{cHead}{RGB}{68,84,106}
\definecolor{cAcc}{RGB}{0,102,153}
\definecolor{cGris}{gray}{0.42}
\definecolor{cSuave}{RGB}{240,244,248}
\definecolor{cBien}{RGB}{34,124,68}
\definecolor{cMal}{RGB}{176,42,42}
\definecolor{cBienSuave}{RGB}{230,243,234}
\definecolor{cMalSuave}{RGB}{250,235,233}
\definecolor{cAccSuave}{RGB}{225,238,246}

\tikzset{
  caja/.style     = {rectangle, rounded corners=2.5pt, draw=cHead, line width=.7pt,
                     fill=cSuave, align=center, inner sep=6pt, font=\small},
  cajaAcc/.style  = {caja, draw=cAcc, fill=cAccSuave},
  cajaOk/.style   = {caja, draw=cBien, fill=cBienSuave},
  cajaNo/.style   = {caja, draw=cMal, fill=cMalSuave, dashed},
  cajaTenue/.style= {caja, draw=cGris, fill=white},
  rotulo/.style   = {font=\scriptsize\itshape, text=cGris, align=center},
  etiq/.style     = {font=\scriptsize, text=cHead, align=center,
                     fill=white, inner sep=1.5pt},
  fl/.style       = {-{Stealth[length=4.5pt]}, draw=cHead, line width=.7pt},
  flAcc/.style    = {fl, draw=cAcc},
  flNo/.style     = {-{Stealth[length=4.5pt]}, draw=cMal, line width=.7pt, dashed},
  llave/.style    = {decorate, decoration={brace, amplitude=4pt, raise=2pt},
                     draw=cGris, line width=.6pt},
  titulo/.style   = {font=\small\bfseries, text=cHead},
}

% Cabecera de figura: #1 número, #2 título
\newcommand{\figtit}[2]{\node[font=\footnotesize\bfseries, text=cGris, anchor=west]}

\begin{document}
\begin{tikzpicture}[node distance=0pt]

\def\an{2.75}   % ancho de caja
\def\se{0.42}   % separación

% --- las cuatro posiciones de la escala ----------------------------
\foreach \i/\tit/\sub [count=\n from 0] in {
  1/Identificable/{nombre, n.\textordmasculine{} de historia,\\fecha exacta},
  2/{Seudonimizado\\clave accesible}/{el equipo puede\\reidentificar},
  3/{Seudonimizado\\clave separada}/{la custodia alguien\\ajeno al análisis},
  4/Anonimizado/{no existe clave;\\irreversible}}
{
  \pgfmathsetmacro{\x}{\n*(\an+\se)}
  \node[caja, text width=\an cm-12pt, minimum height=1.15cm] (c\i) at (\x,0) {\tit};
  \node[rotulo, below=3pt of c\i.south, text width=\an cm] {\sub};
}

% --- flechas entre cajas -------------------------------------------
\foreach \a/\b in {c1/c2, c2/c3, c3/c4}
  \draw[fl] (\a) -- (\b);

% --- frontera del RGPD ---------------------------------------------
\coordinate (fr) at ($(c3.east)!0.5!(c4.west)$);
\draw[cMal, line width=1pt, dash pattern=on 3pt off 2pt]
      ($(fr)+(0,1.45)$) -- ($(fr)+(0,-1.55)$);

% --- llaves y rótulos de ámbito ------------------------------------
\draw[llave] ($(c1.north west)+(0,.12)$) -- ($(c3.north east)+(0,.12)$)
  node[midway, above=6pt, font=\small\bfseries, text=cMal]
  {Datos personales · se aplica el RGPD};

\draw[llave] ($(c4.north west)+(0,.12)$) -- ($(c4.north east)+(0,.12)$)
  node[midway, above=6pt, font=\small\bfseries, text=cBien]
  {Fuera del RGPD};

% --- gradiente de riesgo -------------------------------------------
\draw[{Stealth[length=5pt]}-{Stealth[length=5pt]}, draw=cGris, line width=.6pt]
  ($(c1.south west)+(0,-1.15)$) -- ($(c4.south east)+(0,-1.15)$);
\node[rotulo, anchor=west] at ($(c1.south west)+(.05,-1.42)$) {mayor riesgo de reidentificación};
\node[rotulo, anchor=east] at ($(c4.south east)+(-.05,-1.42)$) {menor riesgo};

% --- aviso sobre la anonimización real ------------------------------
\node[rotulo, text width=5.4cm, anchor=north, text=cMal]
  at ($(c4.south)+(0,-1.95)$)
  {En investigación biomédica la anonimización\\ real es excepcional: casi todo lo que se\\ llama «anonimizado» está en la casilla anterior.};

\end{tikzpicture}
\end{document}
